% GENERATED FILE. Edit canonical lessons, then run npm run generate:books.
% canonical-source-hash: fnv1a64:ca4bb0152467b996
% canonical-lessons: SA-C365-siddhantah, SA-C365-tattvam, SA-C365-prajna, SA-C365-jijnasa, SA-C365-upalaksayati

\chapter{Theory and Practice}
\label{ch:sa-theory-and-practice}
\hlchaptermodality{365}

\begin{chapteropening}
\textbf{By the end of this chapter:} \emph{I can say whether an idea holds up in practice.}

Complete the last lesson of chapter 365: I can say whether an idea holds up in practice.
\end{chapteropening}

\section[\texorpdfstring{\sk{सिद्धान्तः} (siddhāntaḥ)}{सिद्धान्तः (siddhāntaḥ)}]{\sk{सिद्धान्तः} (siddhāntaḥ) — a principle; a theory}

\label{lesson:SA-C365-siddhantah}



\begin{quote}
Before the new one: say the Sanskrit for an essay, then the Sanskrit for investigates.
\end{quote}

\subsection*{You'll want to know: \sk{सिद्धान्तः}}
\textbf{\sk{सिद्धान्तः}} — \emph{siddhāntaḥ} — \textquotedblleft{}a principle; a theory\textquotedblright{}.

\begin{quote}
\sk{सिद्धान्तः} \sk{सरलः}, \sk{किन्तु} \sk{प्रयोगः} \sk{कठिनः।}
\end{quote}

\begin{quote}
\emph{siddhāntaḥ saralaḥ, kintu prayogaḥ kaṭhinaḥ.}
\end{quote}

\begin{quote}
The principle is simple, but putting it into practice is hard.
\end{quote}

A masculine noun.

\begin{cousinweb}
\textbf{\sk{सिद्धान्तः}} (\emph{siddhāntaḥ}) is built from \emph{siddha}, \textquotedblleft{}established\textquotedblright{}, and \emph{anta}, \textquotedblleft{}an end\textquotedblright{}; MW gives \emph{siddhānta} \textquotedblleft{}an established end\textquotedblright{} and \textquotedblleft{}a demonstrated conclusion\textquotedblright{} (MW).
\end{cousinweb}

\begin{grammarlens}[title={What you've built}]
One more word for giving reasons and opinions.
\end{grammarlens}

\subsection*{Your turn}
\begin{itemize}
\raggedright
  \item \emph{Say it:} \emph{siddhāntaḥ}
  \item \emph{Say it:} \emph{siddhāntaḥ}, once more
  \item \emph{Say it:} \emph{siddhāntaḥ saralaḥ, kintu prayogaḥ kaṭhinaḥ.} — The principle is simple, but putting it into practice is hard.
  \item \emph{Recall:} say the Sanskrit for an essay, then the Sanskrit for investigates, then say \emph{siddhāntaḥ} again
\end{itemize}

\begin{tcolorbox}[breakable,colback=teal!4,colframe=teal!35!black,title={Before you move on}]
What does \sk{सिद्धान्तः} mean? (\textquotedblleft{}A principle; a theory\textquotedblright{}.) Say it once more.
\end{tcolorbox}
\section[\texorpdfstring{\sk{तत्त्वम्} (tattvam)}{तत्त्वम् (tattvam)}]{\sk{तत्त्वम्} (tattvam) — reality; the truth of a thing}

\label{lesson:SA-C365-tattvam}



\begin{quote}
Before the new one: say the Sanskrit for investigates, then the Sanskrit for a principle.
\end{quote}

\subsection*{You'll want to know: \sk{तत्त्वम्}}
\textbf{\sk{तत्त्वम्}} — \emph{tattvam} — \textquotedblleft{}reality; the truth of a thing\textquotedblright{}.

\begin{quote}
\sk{पण्डितः} \sk{धर्मस्य} \sk{तत्त्वं} \sk{जानाति।}
\end{quote}

\begin{quote}
\emph{paṇḍitaḥ dharmasya tattvaṃ jānāti.}
\end{quote}

\begin{quote}
The scholar knows the true nature of dharma.
\end{quote}

A neuter noun.

\begin{cousinweb}
\textbf{\sk{तत्त्वम्}} (\emph{tattvam}) is built from \emph{tat}, \textquotedblleft{}that\textquotedblright{}, with the abstract ending \emph{-tva}; MW gives \emph{tat-tva} \textquotedblleft{}true or real state, truth, reality\textquotedblright{} (MW).
\end{cousinweb}

\begin{grammarlens}[title={What you've built}]
One more word for giving reasons and opinions.
\end{grammarlens}

\subsection*{Your turn}
\begin{itemize}
\raggedright
  \item \emph{Say it:} \emph{tattvam}
  \item \emph{Say it:} \emph{tattvam}, once more
  \item \emph{Say it:} \emph{paṇḍitaḥ dharmasya tattvaṃ jānāti.} — The scholar knows the true nature of dharma.
  \item \emph{Recall:} say the Sanskrit for investigates, then the Sanskrit for a principle, then say \emph{tattvam} again
\end{itemize}

\begin{tcolorbox}[breakable,colback=teal!4,colframe=teal!35!black,title={Before you move on}]
What does \sk{तत्त्वम्} mean? (\textquotedblleft{}Reality; the truth of a thing\textquotedblright{}.) Say it once more.
\end{tcolorbox}
\section[\texorpdfstring{\sk{प्रज्ञा} (prajñā)}{प्रज्ञा (prajñā)}]{\sk{प्रज्ञा} (prajñā) — wisdom}

\label{lesson:SA-C365-prajna}



\begin{quote}
Before the new one: say the Sanskrit for a principle, then the Sanskrit for reality.
\end{quote}

\subsection*{You'll want to know: \sk{प्रज्ञा}}
\textbf{\sk{प्रज्ञा}} — \emph{prajñā} — \textquotedblleft{}wisdom\textquotedblright{}.

\begin{quote}
\sk{वृद्धस्य} \sk{प्रज्ञा} \sk{महती} \sk{भवति।}
\end{quote}

\begin{quote}
\emph{vṛddhasya prajñā mahatī bhavati.}
\end{quote}

\begin{quote}
An elder's wisdom is great.
\end{quote}

A feminine noun.

\begin{cousinweb}
\textbf{\sk{प्रज्ञा}} (\emph{prajñā}) is built on \emph{pra} and the root √jñā, \textquotedblleft{}to know\textquotedblright{}; MW gives \emph{pra-jñā} \textquotedblleft{}wisdom, intelligence, knowledge, judgement\textquotedblright{} (MW).
\end{cousinweb}

\begin{grammarlens}[title={What you've built}]
One more word for giving reasons and opinions.
\end{grammarlens}

\subsection*{Your turn}
\begin{itemize}
\raggedright
  \item \emph{Say it:} \emph{prajñā}
  \item \emph{Say it:} \emph{prajñā}, once more
  \item \emph{Say it:} \emph{vṛddhasya prajñā mahatī bhavati.} — An elder's wisdom is great.
  \item \emph{Recall:} say the Sanskrit for a principle, then the Sanskrit for reality, then say \emph{prajñā} again
\end{itemize}

\begin{tcolorbox}[breakable,colback=teal!4,colframe=teal!35!black,title={Before you move on}]
What does \sk{प्रज्ञा} mean? (\textquotedblleft{}Wisdom\textquotedblright{}.) Say it once more.
\end{tcolorbox}
\section[\texorpdfstring{\sk{जिज्ञासा} (jijñāsā)}{जिज्ञासा (jijñāsā)}]{\sk{जिज्ञासा} (jijñāsā) — curiosity; the wish to know}

\label{lesson:SA-C365-jijnasa}



\begin{quote}
Before the new one: say the Sanskrit for reality, then the Sanskrit for wisdom.
\end{quote}

\subsection*{You'll want to know: \sk{जिज्ञासा}}
\textbf{\sk{जिज्ञासा}} — \emph{jijñāsā} — \textquotedblleft{}curiosity; the wish to know\textquotedblright{}.

\begin{quote}
\sk{बालेषु} \sk{बहु} \sk{जिज्ञासा} \sk{भवति।}
\end{quote}

\begin{quote}
\emph{bāleṣu bahu jijñāsā bhavati.}
\end{quote}

\begin{quote}
Children have a great deal of curiosity.
\end{quote}

A feminine noun.

\begin{cousinweb}
\textbf{\sk{जिज्ञासा}} (\emph{jijñāsā}) comes from the root √jñā, \textquotedblleft{}to know\textquotedblright{}, in its desiderative form, \textquotedblleft{}to wish to know\textquotedblright{}; MW gives \emph{jijñāsā} as the desire of knowing (MW).
\end{cousinweb}

\begin{grammarlens}[title={What you've built}]
One more word for giving reasons and opinions.
\end{grammarlens}

\subsection*{Your turn}
\begin{itemize}
\raggedright
  \item \emph{Say it:} \emph{jijñāsā}
  \item \emph{Say it:} \emph{jijñāsā}, once more
  \item \emph{Say it:} \emph{bāleṣu bahu jijñāsā bhavati.} — Children have a great deal of curiosity.
  \item \emph{Recall:} say the Sanskrit for reality, then the Sanskrit for wisdom, then say \emph{jijñāsā} again
\end{itemize}

\begin{tcolorbox}[breakable,colback=teal!4,colframe=teal!35!black,title={Before you move on}]
What does \sk{जिज्ञासा} mean? (\textquotedblleft{}Curiosity; the wish to know\textquotedblright{}.) Say it once more.
\end{tcolorbox}
\section[\texorpdfstring{\sk{उपलक्षयति} (upalakṣayati)}{उपलक्षयति (upalakṣayati)}]{\sk{उपलक्षयति} (upalakṣayati) — notices; observes}

\label{lesson:SA-C365-upalaksayati}



\begin{quote}
Before the new one: say the Sanskrit for wisdom, then the Sanskrit for curiosity.
\end{quote}

\subsection*{You'll want to know: \sk{उपलक्षयति}}
\textbf{\sk{उपलक्षयति}} — \emph{upalakṣayati} — \textquotedblleft{}notices; observes\textquotedblright{}.

\begin{quote}
\sk{वैज्ञानिकः} \sk{नक्षत्राणि} \sk{उपलक्षयति।}
\end{quote}

\begin{quote}
\emph{vaijñānikaḥ nakṣatrāṇi upalakṣayati.}
\end{quote}

\begin{quote}
The scientist observes the stars.
\end{quote}

Past (he …): \sk{सः} \sk{दोषम्} \sk{उपलक्षितवान्}, \emph{saḥ doṣam upalakṣitavān}.

\begin{cousinweb}
\textbf{\sk{उपलक्षयति}} (\emph{upalakṣayati}) is built on \emph{upa}, \textquotedblleft{}towards\textquotedblright{}, and the verb \emph{lakṣaya}, \textquotedblleft{}to mark\textquotedblright{}; MW gives \emph{upa-√lakṣ} \textquotedblleft{}to observe, to perceive\textquotedblright{} and \textquotedblleft{}to pay attention to\textquotedblright{} (MW).
\end{cousinweb}

\begin{grammarlens}[title={What you've built}]
One more word for giving reasons and opinions.
\end{grammarlens}

\subsection*{Your turn}
\begin{itemize}
\raggedright
  \item \emph{Say it:} \emph{upalakṣayati}
  \item \emph{Say it:} \emph{upalakṣayati}, once more
  \item \emph{Say it:} \emph{vaijñānikaḥ nakṣatrāṇi upalakṣayati.} — The scientist observes the stars.
  \item \emph{Recall:} say the Sanskrit for wisdom, then the Sanskrit for curiosity, then say \emph{upalakṣayati} again
\end{itemize}

\begin{tcolorbox}[breakable,colback=teal!4,colframe=teal!35!black,title={Before you move on}]
What does \sk{उपलक्षयति} mean? (\textquotedblleft{}Notices; observes\textquotedblright{}.) Say it once more.
\end{tcolorbox}
